\documentclass{situationist-preprint}

\title{The Problem of Substitutes}
\msid{EXI-16-b}
\msfamily{EXI-16}
\msdiscipline{Meaning, Death and the Absurd}
\mstype{Companion manuscript}
\mskeywords{love; substitutes; trading up; duplication; relational reasons; irreplaceability}

\begin{document}
\maketitle

\begin{abstract}
If love responds to reasons, a substitute with the same or better qualities should be as good an object of love as the beloved. EXI-16-a answers that established love answers to relational reasons, which no qualitatively similar person shares. We extend the answer to three cases of substitution: trading up to a person with better qualities, replacement by a perfect duplicate, and the new love that follows a loss. In each, the substitute lacks the relational reason, and in the third the new love answers to a relational reason of its own rather than inheriting the old one. The same structure underlies the non-transferability of partial reasons in ETH-19-b.
\end{abstract}

\section{Three cases}

The problem of substitutes is old \citep[208--212]{PlatoSymp}, and its modern form is due to the observation that if we love people for their qualities, then whoever has the qualities should do as well \citep{Nozick1989,Kolodny2003}. EXI-16-a distinguishes qualitative reasons, features of the beloved that do not involve the lover, from relational reasons, a shared history between them (Definition~2.2 of EXI-16-a), and shows that qualitative reasons alone would make love transferable (Lemma~3.2) while relational reasons block the transfer (Lemma~3.3). Since every established love answers to relational reasons (Lemma~3.4), established love is not transferable (Theorem~3.5). We apply the result to three cases.

\begin{proposition}[Trading up]\label{p:up}
That another person has the qualities for which $L$ loves $x$ to a higher degree is not a reason for $L$ to love that person instead of $x$ that is at least as strong as the reasons for $L$'s love of $x$.
\end{proposition}

\begin{proof}
By Lemma~3.4 of EXI-16-a, $L$'s established love of $x$ answers to relational reasons. By Principle~P3 of EXI-16-a, the other person does not share the history that constitutes them. So the reasons for loving her include only qualitative ones, and whatever they gain in degree, they lack a kind of reason that $L$'s love of $x$ has. The reasons for $L$'s love of $x$ therefore do not favor the other person at least as strongly.
\end{proof}

\begin{proposition}[Duplication]\label{p:dup}
If the beloved is duplicated with her memories, either the duplicate is the beloved, and the love continues, or she is not, and she lacks the relational reasons for which the beloved is loved.
\end{proposition}

\begin{proof}
This is the reply to the duplication objection in EXI-16-a. Apparent memories of a shared history are qualitative features of the one who has them; having taken part in the history is not. Whether the duplicate took part in it depends on whether she is the beloved, which is settled by the correct account of personal identity (PID-01), not by the theory of love.
\end{proof}

\begin{proposition}[After loss]\label{p:after}
A new love that follows the loss of a beloved is not the old love transferred. It answers to a relational reason of its own, constituted by a new history.
\end{proposition}

\begin{proof}
If the new love answered to the reasons of the old, it would answer to relational reasons constituted by the history with the lost beloved; but by Principle~P3 of EXI-16-a, the new beloved did not take part in that history. So by Lemma~3.4 of EXI-16-a, once it is established, the new love answers to a relational reason constituted by the history between $L$ and the new beloved, which is a different reason.
\end{proof}

Proposition~\ref{p:after} explains the fittingness of grief. If love were transferable, the loss of the beloved would be the loss of an instance of qualities that could be found again, and grief would be fitting only to the extent that substitutes are scarce. Since the relational reason is lost with the person, the loss is of something that no new love restores, and grief is fitting even when a new love comes.

\section{The same structure elsewhere}

ETH-19-b argues that the reason that suffices for a man to save his wife, that she is his wife, is relational and non-transferable, while an impartial justification would supply a reason that favors anyone standing in a relationship with the same welfare effects (Proposition~2.1 of ETH-19-b). This is the structure of Lemmas~3.2 and 3.3 of EXI-16-a applied to reasons for action: a reason that consists in features anyone could have generalizes to anyone who has them, by Principle~P2 of EXI-16-a, and a reason that consists in a shared history does not. The problem of substitutes and the problem of the thought too many have one solution.

\begin{objection}[Velleman]
\citet{Velleman1999} holds that love responds to the value of the beloved's rational nature, which every person has, and that its object is irreplaceable because that value is a dignity, which admits no substitutes.
\end{objection}

\begin{reply}
Velleman's view explains why no person may be traded for another, but it does not explain why $L$ loves this person rather than any other, since the value it cites is shared by all. Relational reasons supply what is missing: they explain the selectivity of love without denying that its object has a dignity that admits no price.
\end{reply}

\section{Conclusion}

A substitute, however similar or however much better, is not loved for the reasons for which the beloved is loved. The reason it lacks is the shared history, and a new love, if one comes, answers to a history of its own.

\begin{thebibliography}{9}\small
\bibitem[Kolodny(2003)]{Kolodny2003} Kolodny, N. (2003). Love as valuing a relationship. \emph{The Philosophical Review}, 112(2), 135--189.
\bibitem[Nozick(1989)]{Nozick1989} Nozick, R. (1989). Love's bond. In \emph{The Examined Life: Philosophical Meditations}. Simon and Schuster.
\bibitem[Plato(Symp.)]{PlatoSymp} Plato. \emph{Symposium}. Cited by Stephanus pages.
\bibitem[Velleman(1999)]{Velleman1999} Velleman, J.~D. (1999). Love as a moral emotion. \emph{Ethics}, 109(2), 338--374.
\end{thebibliography}

\end{document}
